\section{Introduction}
\label{sec:intro}

A networked game has to decide where the rules live. If the client enforces
them, every client is part of the trusted computing base and a single modified
or buggy one can corrupt the game. If the server enforces them, the client can
be as simple and as untrusted as you like. This project takes the second
position and pushes it to its limit, as the brief requires~\cite{brief}: the
microcontroller is the only component that may change state, and the PC
application is a view with a mouse attached. The property the design guarantees
is narrow and testable, namely that a rejected request never changes the game
state.

\textbf{The game.} Tic-tac-toe was chosen because its rules are small enough to
implement exactly and to test exhaustively, which puts the interesting work in
the protocol and the authority model rather than in game logic. A nine-cell
board also encodes in nine ASCII characters, so a full board fits in one short
line and the whole state can be sent after every move instead of a delta. That
removes an entire class of client and server disagreeing about state.

\textbf{The GUI library.} Tkinter~\cite{tkinter} was chosen because it ships
with CPython, so the marker needs no package installation to run the client, and
because a three-by-three grid of buttons is exactly what it is good at. It also
has a timer facility that makes it straightforward to keep the window responsive
while a reply is outstanding, which Section~\ref{sec:design} describes. A
heavier toolkit would have added a dependency and bought nothing the application
needs.

\textbf{Design choices.} The rules and the frame parser are plain C with no
hardware abstraction layer (HAL) and no hardware headers, so the same sources
compile for the STM32 and for a host simulator. One file isolates everything
board-specific behind a handful of functions. Messages are framed as single
newline-terminated lines of printable ASCII with a bounded length, which makes
the parser a small finite-state machine (FSM) that can be reasoned about and
fuzzed from a test. The practical consequence is that a rule bug is found in
milliseconds on a laptop rather than on the bench with a debugger.

\textbf{Task allocation.} This project was done individually, so there is no
split of work between two students to report. It was nonetheless organised as if
it were a team: each piece of functionality was developed on its own branch and
merged through a pull request, with the game core, the protocol, the GUI, the
board application and the documentation each landing separately. Reviewing one's
own pull request catches less than a second reader would, but it still forces
the diff to be read once in isolation, keeps unrelated changes from travelling
together, and leaves a history in which each change can be read, and if
necessary reverted, on its own.

\textbf{Structure.} Section~\ref{sec:protocol} defines the message protocol,
since every other part of the system is specified against it.
Section~\ref{sec:design} gives the responsibility split, the server FSM, the
client flow and the hardware connection. Section~\ref{sec:impl} describes how
the firmware is put together, Section~\ref{sec:authority} argues the authority
property from the structure of the code, and Section~\ref{sec:robust} covers
what the parser does with hostile input. Section~\ref{sec:valid} reports the
automated validation and Section~\ref{sec:conclusion} draws conclusions and
names the limitations.
